\documentclass[11pt]{article}
% Shared design for the two minimal papers.
% It follows the layout of the guide "The Riemann Hypothesis and its generalizations"
% (KNT with Fable, 9 November 2026): a designed title block, a short story list at the start,
% coloured headings, boxed statements, idea boxes, and annotated figures.
%
% A paper sets \papertitle, \papersubtitle and \paperrunning before \begin{document}
% and calls \maketitleblock after it.

\usepackage[letterpaper,left=1.12in,right=1.12in,top=1.05in,bottom=1.1in]{geometry}
\usepackage{amsmath,mathtools,amssymb,amsthm}
\usepackage{newpxtext}
\usepackage[varbb]{newpxmath}
\usepackage{fontspec}
\setsansfont{texgyreheros}[Extension=.otf, UprightFont=*-regular, BoldFont=*-bold,
  ItalicFont=*-italic, BoldItalicFont=*-bolditalic, Scale=0.88]
\usepackage{microtype}
\usepackage{parskip}
\usepackage{xcolor}
\usepackage{booktabs,array}
\usepackage{enumitem}
\usepackage{multicol}
\usepackage[titles]{tocloft}
\usepackage{tikz}
\usetikzlibrary{decorations.markings,arrows.meta,positioning,calc}
\usepackage[most]{tcolorbox}
\usepackage{fancyhdr}
\usepackage{titlesec}
\usepackage[font=small,labelfont={bf,sf},labelsep=period,skip=6pt]{caption}
\usepackage{url}

\definecolor{accent}{HTML}{1B4F72}
\definecolor{warm}{HTML}{C0562A}
\definecolor{slate}{HTML}{56636E}
\definecolor{ink}{HTML}{1E1E1E}

\usepackage[colorlinks=true,linkcolor=accent,citecolor=accent,urlcolor=accent]{hyperref}

% Lean identifiers: typewriter, breakable at dots and underscores.
\DeclareUrlCommand\lean{\urlstyle{tt}}

% ---------- layout ----------
\linespread{1.06}
\setlength{\parskip}{0.55em plus 0.1em minus 0.05em}
\setlist{itemsep=0.25em, topsep=0.2em, parsep=0.1em}
\titleformat{\section}{\sffamily\bfseries\Large\color{accent}}{\thesection}{0.7em}{}
\titleformat{\subsection}{\sffamily\bfseries\normalsize\color{ink}}{\thesubsection}{0.6em}{}
\titleformat{\paragraph}[runin]{\sffamily\bfseries\small\color{accent}}{}{0pt}{}[\enspace]
\titlespacing*{\section}{0pt}{22pt}{5pt}
\titlespacing*{\subsection}{0pt}{13pt}{1pt}
\titlespacing*{\paragraph}{0pt}{9pt}{0pt}
\pagestyle{fancy}
\fancyhf{}
\renewcommand{\headrulewidth}{0pt}
\fancyhead[L]{\sffamily\footnotesize\color{slate}\paperrunning}
\fancyhead[R]{\sffamily\footnotesize\color{slate}\thepage}
\setlength{\headheight}{14pt}
\setcounter{tocdepth}{1}
\setlength{\cftbeforesecskip}{2pt}
\renewcommand{\cftsecfont}{\small}
\renewcommand{\cftsecpagefont}{\small\color{slate}}
\renewcommand{\cftsecleader}{\cftdotfill{\cftdotsep}}
\setlength{\cftsecnumwidth}{1.7em}
\setlength{\emergencystretch}{2em}
\numberwithin{equation}{section}

% ---------- title block ----------
\newcommand{\papertitle}{}
\newcommand{\papersubtitle}{}
\newcommand{\paperrunning}{}
\newcommand{\paperauthors}{KNT\enspace\textperiodcentered\enspace Astra\enspace\textperiodcentered\enspace Fable~5.2}
\newcommand{\paperdate}{Draft of 11 November 2026}
\newcommand{\maketitleblock}{%
  \thispagestyle{plain}%
  {\sffamily\bfseries\color{accent}\fontsize{25}{29}\selectfont\raggedright\papertitle\par}%
  \vspace{5pt}%
  {\sffamily\Large\color{slate}\raggedright\papersubtitle\par}%
  \vspace{8pt}%
  {\color{accent}\hrule height 1.2pt}%
  \vspace{7pt}%
  {\sffamily\small\color{ink}\paperauthors\hfill\color{slate}\paperdate\par}%
  \vspace{10pt}}

% Small headings inside the front matter.
\newcommand{\fronthead}[1]{\par\vspace{6pt}{\sffamily\bfseries\normalsize\color{ink}#1}\par\vspace{-2pt}}
\makeatletter
\newcommand{\contentsblock}{%
  \par\noindent\begin{minipage}{\linewidth}%
  \fronthead{Contents}%
  \setlength{\parskip}{0pt}\setlength{\columnsep}{2.2em}\raggedright
  \begin{multicols}{2}
  \@starttoc{toc}
  \end{multicols}
  \end{minipage}\par}
\makeatother

% ---------- statements ----------
% Theorems, lemmas, propositions and corollaries share one counter.
% Each statement is in a tinted box with an accent rule at the left.
\newtheoremstyle{boxed}{0pt}{0pt}{\normalfont}{0pt}{\sffamily\bfseries\small\color{accent}}{.}{0.6em}{}
\theoremstyle{boxed}
\newtheorem{theorem}{Theorem}[section]
\newtheorem{lemma}[theorem]{Lemma}
\newtheorem{proposition}[theorem]{Proposition}
\newtheorem{corollary}[theorem]{Corollary}

\tcbset{guidebox/.style={enhanced, parbox=false, boxrule=0pt, frame hidden,
  arc=1.5pt, left=11pt, right=10pt, top=7pt, bottom=7pt, before skip=11pt, after skip=9pt}}
% Not breakable: a label at the start of a box that breaks at the top of a page can be lost.
\tcbset{statementbox/.style={guidebox, colback=accent!6,
  borderline west={3pt}{0pt}{accent}}}
\tcolorboxenvironment{theorem}{statementbox}
\tcolorboxenvironment{lemma}{statementbox}
\tcolorboxenvironment{proposition}{statementbox}
\tcolorboxenvironment{corollary}{statementbox}

% The main result of a paper: a framed box with a title bar.
% \begin{maintheorem}{name} uses the counter of the theorems, and \label works after it.
\newtcolorbox{mainresult}[1]{enhanced, parbox=false, colback=white, colframe=accent, boxrule=0.9pt,
  arc=1.5pt, left=11pt, right=11pt, top=8pt, bottom=8pt, title={#1},
  fonttitle=\sffamily\bfseries\small, colbacktitle=accent, coltitle=white,
  before skip=12pt, after skip=12pt}
\newenvironment{maintheorem}[1]%
  {\refstepcounter{theorem}\begin{mainresult}{Theorem~\thetheorem\ (#1)}}%
  {\end{mainresult}}

% ---------- proofs ----------
% A proof has a thin rule at the left, so that the eye finds the statements.
\renewcommand{\proofname}{\upshape\sffamily\bfseries\small\color{slate}Proof}
\renewcommand{\qedsymbol}{\textcolor{slate}{\rule{0.55em}{0.55em}}}
\tcolorboxenvironment{proof}{enhanced, breakable, parbox=false, boxrule=0pt, frame hidden,
  colback=white, borderline west={1.2pt}{0pt}{slate!35}, left=11pt, right=0pt, top=1pt, bottom=2pt,
  before skip=6pt, after skip=11pt}
% A step or a case inside a proof.
\newcommand{\step}[1]{{\sffamily\bfseries\small\color{slate}#1}}

% ---------- idea, notation and fine-print boxes ----------
\newtcolorbox{idea}[1]{guidebox, breakable, colback=warm!6, borderline west={3pt}{0pt}{warm!80},
  before upper={{\sffamily\bfseries\footnotesize\color{warm!75!black}IDEA\enspace\textbar\enspace}{\sffamily\bfseries\small #1}\par}}
\newtcolorbox{notation}[1]{guidebox, breakable, colback=black!4, borderline west={3pt}{0pt}{slate!60},
  fontupper=\small,
  before upper={{\sffamily\bfseries\footnotesize\color{slate}NOTATION\enspace\textbar\enspace}{\sffamily\bfseries\small #1}\par}}
\newtcolorbox{fineprint}{guidebox, breakable, colback=white, borderline west={1.5pt}{0pt}{black!25},
  fontupper=\footnotesize\color{black!78}, left=9pt, top=1pt, bottom=1pt,
  before upper={{\sffamily\bfseries Fine print.}\enspace}}

% ---------- figures ----------
\tikzset{
  vtx/.style={circle, fill=ink, inner sep=1.35pt},
  port/.style={circle, draw=accent, fill=white, line width=0.8pt, inner sep=1.5pt},
  edge/.style={line width=0.55pt, draw=ink},
  hot/.style={line width=1.1pt, draw=warm},
  cool/.style={line width=1.1pt, draw=accent},
  region/.style={draw=accent, fill=accent!7, line width=0.6pt, rounded corners=6pt},
  note/.style={font=\sffamily\footnotesize, text=slate},
}

% ---------- macros of the mathematics ----------
\newcommand{\Z}{\mathcal{Z}}
\newcommand{\PP}{\mathbb{P}}
\newcommand{\EE}{\mathbb{E}}
\newcommand{\comp}{\operatorname{comp}}
\newcommand{\tr}{\operatorname{tr}}
\newcommand{\lgk}[1]{\log^{(#1)}}
\newcommand{\lgs}{\log^{*}}
\newcommand{\sm}{\setminus}
\newcommand{\half}{\tfrac12}
\newcommand{\cO}{\mathcal{O}}


\renewcommand{\papertitle}{Slabs:\\ pieces at girth five without the lemma on rigid sets}
\renewcommand{\papersubtitle}{A fifth research note on ``The maximum number of cycles in a graph of large rank''}
\renewcommand{\paperrunning}{Slabs}
\renewcommand{\paperauthors}{Fable~5.2, for KNT}
\renewcommand{\paperdate}{Research note of 13 November 2026}
\hypersetup{pdftitle={Slabs: pieces at girth five without the lemma on rigid sets},pdfauthor={Fable 5.2, for KNT}}

\newcommand{\inP}[1]{[P,~#1]}

\begin{document}
\maketitleblock

On 12 November 2026 I wrote a formal proof of \inP{Theorem~7.1} at girth five in Lean. Here [P]
is the paper of 36~pages (commit \texttt{805e2ba}). The formal proof of \inP{Lemma~7.6} (pieces at
girth five) has an inner step that is not in the paper. KNT asked for a draft of this step as
text, and for a measurement. This note has the draft in two forms, with complete proofs, and the
numbers. The manuscript is not changed.

\fronthead{The result in five sentences}
\begin{enumerate}[leftmargin=1.6em]
\item In the proof of \inP{Lemma~7.6} a part $M$ of the piece is long and thin, and its boundary
vertices are at its two ends. So $M$ has many disjoint slabs across a shortest path, and for one
slab $B$ every vertex of a given set $U$ is on one of the two sides of~$B$.
\item For this slab a $U$-join of $M$ meets $M[B]$ in an even subgraph or in a join of two
vertices. So a rigid set $U$ gives a full pair in $M[B]$, and \inP{Lemma~7.4} applies. This
replaces \inP{Lemma~7.5} (rigid sets) with its tree, its segments and its count.
\item Form \emph{slabs}: \inP{Lemma~7.5} is removed and \inP{Lemma~7.4} stays. Measured in a copy:
1288 typeset lines in the place of 1305, and 36 statements in the place of 37; 36 pages and 15
figures as before.
\item Form \emph{slabs and rails}: in addition \inP{Lemma~7.4} gets a proof without ladders, by a
count beside a shortest path and a fact on seven integers. Section~7 then uses no result of
Section~5 of the paper, and the text of Section~7.2 is the formal proof. Measured: 1308 lines,
so three lines more than now.
\item The formal proof in Lean has the second form, with other constants. The first form
differs from it in one proof, that of \inP{Lemma~7.4}. I recommend the form \emph{slabs} if the
aim is fewer steps, and the form \emph{slabs and rails} if the aim is a text that the formal
proof follows. The gain is small in both forms.
\end{enumerate}

\section{The measurement}

The script \texttt{research/apply\_slabs.py} changes a copy of the folder
\texttt{exact-maximum}; with the argument \texttt{rails} it applies the second form. The copies
are \texttt{research/trial-slabs.pdf} and \texttt{research/trial-slabs-rails.pdf}. Both build
without an overfull box and without an undefined reference. A typeset line is a line of
\texttt{pdftotext -layout} that is not empty and is not a running head.

\begin{center}\small
\begin{tabular}{lrrrrrr}
\toprule
 & pages & lines & Section 7 & Section 7.2 & statements & figures\\
\midrule
the paper now & 36 & 1305 & 378 & 121 & 37 & 15\\
form \emph{slabs} & 36 & 1288 & 361 & 103 & 36 & 15\\
form \emph{slabs and rails} & 36 & 1308 & 377 & 120 & 36 & 15\\
\bottomrule
\end{tabular}
\end{center}

In both forms the figure of the tree through a rigid set is replaced by a figure of the slabs
(Figure~\ref{fig:slabs}). The statements after \inP{Lemma~7.4} get a number that is smaller by
one. No other section changes, with the exception of one sentence on the formal proof in
\inP{Section~1.4}.

What the two forms change in the proof, step by step:

\begin{center}\small
\begin{tabular}{p{0.3\textwidth}p{0.3\textwidth}p{0.3\textwidth}}
\toprule
the paper now & form \emph{slabs} & form \emph{slabs and rails}\\
\midrule
\inP{Lemma~7.4}: a full pair at girth five is small; by the reduction \inP{Lemma~5.2} and the
ladders of \inP{Lemma~5.6} & the same & a rigid pair at girth five is small; by a shortest
path, a forest beside it, and seven integers\\
\addlinespace
\inP{Lemma~7.5}: a thin multigraph with a rigid set is small; a smallest tree through the set,
its segments, a far point, an exchange step, two cases, a count & no such lemma & no such
lemma\\
\addlinespace
\inP{Lemma~7.6}: groups of layers, then \inP{Lemma~7.5} and the product & groups of layers,
slabs, side components, the same two cases, the product & the same as in the form
\emph{slabs}\\
\bottomrule
\end{tabular}
\end{center}

\section{The new proof of Lemma 7.6}

\begin{notation}{what the proof uses from the paper}
\begin{itemize}[leftmargin=1.2em]
\item $N$ is a connected subcubic multigraph. Its vertices of degree less than three are its
boundary vertices, its deficit $d(N)$ is the sum of $3-\deg v$ over all vertices, and
$|V(N)|=2\beta(N)-2+d(N)$; this is \inP{equation~(7.1)}.
\item For an even set $T$ of boundary vertices, a $T$-join is a set of edges in which the
vertices of $T$ have odd degree and all other vertices have even degree. The forest proportion
of $T$ is the proportion of the $T$-joins that are forests, and $T$ is rigid if every $T$-join
is a forest.
\item \inP{Lemma~7.3} (product). Let pairwise disjoint connected sets of vertices of degree
three, without a vertex of $T$, induce multigraphs in which every even set of boundary vertices
has forest proportion at most $f_i$. Then the forest proportion of $T$ in $N$ is at most the
product of the $f_i$.
\item \inP{Lemma~7.4} (full pairs at girth five). Let $M$ be a connected subcubic multigraph of
girth at least five, and let $x\ne y$ be vertices of degree at most two such that every
$\{x,y\}$-join of $M$ is a path. Then $|V(M)|<15\,d(M)$.
\item $B$ is the constant of \inP{Section~7.1}: in a piece with $k$ boundary vertices and a
forest proportion of at least $\frac35$, every distance layer from a boundary vertex has fewer
than $B(k+1)^{2}$ vertices. This is the step ``A marker'' in the proof of \inP{Lemma~7.2}.
\end{itemize}
\end{notation}

\begin{lemma}[Pieces at girth five; this is {\inP{Lemma~7.6}}]\label{lem:pieces}
For every integer $k\ge1$ there is a number $F_0(k)$ with the following property. Let $N$ be a
connected subcubic multigraph of girth at least five with at most $k$ boundary vertices and at
least $F_0(k)$ vertices. Then every set of boundary vertices of even size has forest proportion
less than $\frac35$.
\end{lemma}

\begin{figure}[ht]
\centering
\begin{tikzpicture}[x=1cm,y=1cm]
  \draw[draw=slate!70, fill=black!4, line width=0.6pt, rounded corners=6pt] (0,-1.25) rectangle (13.6,1.25);
  \foreach \a in {1.5,4.2,6.9,9.6} \draw[region] (\a,-0.62) rectangle (\a+1.9,0.95);
  % a branch from the slab B_j back to the first layer
  \draw[edge] (5.0,-0.62) .. controls (4.8,-0.95) and (4.4,-1.0) .. (3.8,-1.0) -- (1.2,-1.0)
    .. controls (0.75,-1.0) and (0.5,-0.85) .. (0.5,-0.55);
  % the path pi
  \draw[hot] (0.5,0) -- (13.1,0);
  \draw[edge] (0.5,0) -- (0.5,0.55) (13.1,0) -- (13.1,0.55) (13.1,0) -- (13.1,-0.55);
  \draw[edge] (7.45,0) -- (7.45,0.6) -- (8.25,0.6) -- (8.25,0);
  \foreach \p in {(5.0,-0.62),(3.8,-1.0),(2.5,-1.0),(1.2,-1.0),(7.45,0.6),(8.25,0.6),(7.45,0),(7.85,0),(8.25,0)}
    \node[vtx] at \p {};
  \foreach \p in {(0.5,0),(13.1,0),(6.55,0),(7.15,0),(8.55,0),(9.15,0)} \node[vtx] at \p {};
  \foreach \p in {(0.5,0.55),(0.5,-0.55),(13.1,0.55),(13.1,-0.55)} \node[port] at \p {};
  \node[note, text=accent] at (7.85,1.5) {the set $B=B_i$};
  \node[note, text=accent] at (5.15,1.5) {$B_j$};
  \node[note, text=accent] at (6.5,-0.3) {$s$};
  \node[note, text=accent] at (7.15,-0.3) {$x$};
  \node[note, text=accent] at (8.55,-0.3) {$y$};
  \node[note, text=accent] at (9.2,-0.3) {$t$};
  \node[note, text=warm] at (3.0,0.28) {$\pi$};
  \node[note, text=accent] at (0.22,-0.55) {$u$};
  \node[note] at (3.3,-1.55) {the component $K^{-}$ of $M-B$};
  \node[note] at (11.4,-1.55) {the component $K^{+}$ of $M-B$};
\end{tikzpicture}
\caption{The proof of Lemma~7.6. The part $M$ with the path $\pi$ (thick), four
of its slabs, and the vertices of $U$ (circles) in the first and in the last layer. The vertex
$u$ is in a side component of $M-B_j$; for the slab $B$ it is in $K^{-}$. If $K^{-}\ne K^{+}$
and each of them contains an odd number of vertices of $U$, then a $U$-join of $M$ meets $M[B]$
in an $\{x,y\}$-join.}
\label{fig:slabs}
\end{figure}



\begin{proof}
Put $W=B(k+1)^{2}$, $\tau=60\,W$, $h=2W\tau+2$ and $F_0(k)=Wh\,(k+2^{Wh}+2)$. Suppose that some
even set $T$ of boundary vertices has forest proportion at least $\frac35$. If $T$ is empty, then
$N$ has no cycle and fewer than $2k$ vertices. In the other case fix a boundary vertex $o$, and
define the layers as in the proof of Lemma~7.2. The step ``A marker'' of that proof
uses no girth, and it shows that each layer has fewer than $W$ vertices. So $N$ has at least
$h\,(k+2^{Wh}+2)$ layers. Divide the layers $1,2,\dots$ into groups of $h$ consecutive layers. At
least $2^{Wh}$ groups contain no boundary vertex. In such a group let $M$ be a component of the
subgraph induced by the group that has a vertex $v$ in the last layer of the group. A shortest
path from $o$ to $v$ has one vertex in each layer of the group; let $\pi_j$ be its vertex in the
$j$th of these layers. The path $\pi=\pi_1\cdots\pi_h$ lies in $M$. The vertices of $M$ have
degree three in $N$, and an edge that leaves $M$ starts in the first or in the last layer of the
group. So the boundary vertices of $M$ are in these two layers, there are fewer than $2W$ of
them, and the rank of $M$ is less than $Wh$. We show that no even set of boundary vertices of
$M$ is rigid (Figure~\ref{fig:slabs}).

\step{Slabs.} For $1\le i\le2W$, the $i$th \emph{slab} consists of the layers
$2+(i-1)\tau,\dots,1+i\tau$ of the group. Let $B_i$ be the vertex set of the component of the
subgraph of $M$ induced by this slab that contains the $\tau$ vertices of $\pi$ in the slab. Let
$s_ix_i$ be the edge of $\pi$ that enters $B_i$ and $y_it_i$ the edge that leaves it, with
$x_i,y_i\in B_i$. The vertices of $B_i$ have degree three in $M$, and an edge of $M$ that leaves
$B_i$ starts in the first or in the last layer of the slab. So $d(M[B_i])<4W\le|B_i|$. Hence
$M[B_i]$ has a cycle, by~(7.1), and $x_i\ne y_i$ are boundary vertices of
$M[B_i]$. Call a component of $M-B_i$ a \emph{side component} if it contains neither $s_i$ nor
$t_i$. The vertices of $\pi$ outside $B_i$ are joined to $s_i$ or to $t_i$ in $\pi-B_i$. A set
$B_j$ with $j\ne i$ is connected, avoids $B_i$ and contains a vertex of $\pi$. So a side
component of $M-B_i$ meets no set $B_j$.

\step{Side components.} Let $K$ be a side component of $M-B_i$, and let $j\ne i$. Some
vertex of $K$ has a neighbour in $B_i$, because $M$ is connected. So $K\cup B_i$ is connected,
avoids $B_j$ and contains a vertex of $\pi$. Hence it lies in the component of $s_j$ or of
$t_j$ in $M-B_j$. So a vertex of $M$ is in a side component of $M-B_i$ for at most one~$i$.

\step{No rigid set.} Let $U$ be an even set of boundary vertices of $M$. It has fewer than $2W$
vertices, and none of them is in a slab. So there is an $i$ such that no vertex of $U$ is in a
side component of $M-B_i$. Write $B$, $s$, $x$, $y$, $t$ for $B_i$, $s_i$, $x_i$, $y_i$, $t_i$,
and let $K^{-}$ and $K^{+}$ be the components of $M-B$ that contain $s$ and $t$. Every vertex
of $U$ is in one of them. Put $U^{-}=U\cap V(K^{-})$. Suppose first that $K^{-}=K^{+}$ or that
$|U^{-}|$ is even. Then $M-B$ contains a $U$-join, and its union with a cycle of $M[B]$ is a
$U$-join that is not a forest. Suppose now that $K^{-}\ne K^{+}$ and that $|U^{-}|$ is odd, and
let $J$ be an $\{x,y\}$-join of $M[B]$. Add the edges $sx$ and $ty$, a join of $K^{-}$ for the
set $U^{-}$ with $s$ added or removed, and a join of $K^{+}$ for the set $U\sm U^{-}$ with $t$
added or removed. The result is a $U$-join of $M$. If $U$ is rigid, then $J$ is a forest, and hence a path. Then
Lemma~7.4 gives $|B|<15\,d(M[B])<60\,W=\tau$, and this is false.

So each even set of boundary vertices of $M$ has forest proportion at most $1-2^{-Wh}$. By
Lemma~7.3, the forest proportion of $T$ is at most
$(1-2^{-Wh})^{2^{Wh}}<\half$. This is a contradiction.
\end{proof}


The proof begins and ends as the proof in the paper does. The steps ``Slabs'', ``Side
components'' and ``No rigid set'' are in the place of the reference to \inP{Lemma~7.5}. The
paper has $h=10^{4}W^{3}$; here $h=120\,W^{2}+2$.

\section{A proof of Lemma 7.4 without ladders}

An $\{x,y\}$-join that is a forest is a path. So the next lemma is \inP{Lemma~7.4} with the
constant $133$ in the place of $15$. With it, the proof of Lemma~\ref{lem:pieces} takes
$\tau=532\,W$, and its last case ends with $|B|<133\,d(M[B])<532\,W=\tau$.

\begin{lemma}[Rigid pairs at girth five]\label{lem:rails}
Let $M$ be a connected subcubic multigraph of girth at least five. Let $x\ne y$ be vertices of
degree at most two such that every $\{x,y\}$-join of $M$ is a forest. Then $|V(M)|<133\,d(M)$.
\end{lemma}

\begin{proof}
Put $d=d(M)$. Let $A=a_0a_1\cdots a_\ell$ be a shortest path from $x$ to $y$, and put
$F=M-V(A)$. The edge set of $A$ is an $\{x,y\}$-join, and so is its union with a cycle of $F$. So
$F$ is a forest. A vertex of $F$ has at most one neighbour on $A$: two neighbours $a_i$ and
$a_j$ with $i<j$ give $j-i\le2$, because $A$ is a shortest path, and then a cycle with at most
four edges. For the same reasons an inner vertex of $A$ of degree three has a neighbour in $F$.

Call a vertex of $F$ \emph{plain} if it has degree three and a neighbour $a_p$ with $0<p<\ell$.
A plain vertex has exactly two neighbours in $F$. A vertex of $F$ that is not plain is a
boundary vertex, or it is adjacent to $x$ or $y$, or it has three neighbours in $F$. Let $n_0$,
$n_1$ and $n_3$ be the numbers of these vertices. Then $n_0\le d-2$ and $n_1\le2$, because $x$
and $y$ are boundary vertices. The forest $F$ has at most as many edges as vertices, so a count
of edge ends gives $3n_3+2(|F|-n_0-n_3)\le2|F|$, and $n_3\le2n_0$. So fewer than $3d$ vertices
of $F$ are not plain.

\step{No path of seven plain vertices.} Suppose that $v_0v_1\cdots v_6$ is a path of plain
vertices in $F$, and let $a_{p_k}$ be the neighbour of $v_k$ on $A$. Let $0\le k<l\le6$. Then
$p_k\ne p_l$, because $a_{p_k}$ has three neighbours only. Also $|p_k-p_l|\le l-k+2$, because
$A$ is a shortest path and $a_{p_k}v_k\cdots v_la_{p_l}$ is a path. And $|p_k-p_{k+1}|\ge2$,
because $M$ has no cycle with four edges. Put $s_k=p_{k+1}-p_k$ for $0\le k\le5$. Then $|s_k|$
is $2$ or~$3$, and $0<|s_k+s_{k+1}|\le4$. So two consecutive steps of the same sign are both
$2$ or both $-2$, and of two consecutive steps of different signs one has absolute value $2$
and the other has absolute value~$3$. No three consecutive steps have the same sign, because
$|p_{k+3}-p_k|\le5$.

If consecutive steps always have different signs, then $s_0=s_2=s_4$ and $s_1=s_3=s_5$. Then
$p_5-p_0=3s_0+2s_1=0$ if $|s_0|=2$, and $p_6-p_1=2s_0+3s_1=0$ if $|s_0|=3$. In the other case
let $s_k=s_{k+1}=2\varepsilon$ with $\varepsilon=\pm1$. Where these steps exist,
$s_{k-1}=s_{k+2}=-3\varepsilon$ and $s_{k-2}=s_{k+3}=2\varepsilon$. Then $p_{k-1}=p_{k+4}$ and
$p_{k-2}=p_{k+3}$, so in each of these two pairs one index is outside the range $0,\dots,6$.
This leaves $k=0$ and $k=4$. Let $k=0$; the case $k=4$ is the same with the order reversed. The
first four steps are $2\varepsilon,2\varepsilon,-3\varepsilon,2\varepsilon$. If
$s_4=-3\varepsilon$, then $p_5=p_0$. If $s_4=2\varepsilon$, then $s_5=-3\varepsilon$ and
$p_6=p_1$. So $F$ has no path of seven plain vertices.

\step{The count.} Let $w$ be a vertex of $F$ such that all vertices at distance at most three
from $w$ in $F$ are plain. Then $w$, its two neighbours in $F$, their other neighbours and the
next two vertices form a path of seven plain vertices, because $F$ is a forest. So every vertex
of $F$ is within distance three in $F$ of a vertex that is not plain, and $|F|<22\cdot3d$. An inner vertex of $A$ of
degree three has a neighbour in $F$, and a vertex of $F$ has at most one neighbour on $A$. So
$A$ has at most $|F|+d$ vertices, and $|V(M)|\le2|F|+d<133\,d$.
\end{proof}



The bound of seven is the best one for this argument: the integers $5,2,4,1,3,0$ satisfy the
three conditions on the positions. The script \texttt{checks/rail\_check.py} tests all
sequences of steps (exit code~0).

\section{The formal proof}

The files are in \texttt{Erdos1016/ExactMaximum/GirthFive/} (commit \texttt{b7f9fc9}); the README
of that folder has the same proofs in words.

\begin{center}\small
\begin{tabular}{p{0.34\textwidth}p{0.6\textwidth}}
\toprule
text of this note & Lean declaration\\
\midrule
the fact on seven integers & \lean{GirthFive.rail_bound}\\
a vertex beside a shortest path has one neighbour on it; no path of seven plain vertices &
\lean{GirthFive.attach_unique}, \lean{GirthFive.no_rail}\\
Lemma~\ref{lem:rails} & \lean{GirthFive.two_terminal_card_le}, \lean{GirthFive.two_terminal_region},
with the bound $400\,(|D|+2)$ for the set $D$ of boundary vertices\\
the steps ``Side components'' and ``No rigid set'' & \lean{GirthFive.side_of_other_slab},
\lean{GirthFive.exists_word_two_terminals}, \lean{GirthFive.region_not_rigid}\\
the step ``Slabs'' & \lean{GirthFive.window_not_rigid}\\
Lemma~\ref{lem:pieces} & \lean{GirthFive.exists_affine_bound_girth_five}\\
\bottomrule
\end{tabular}
\end{center}

Two differences. The formal step ``No rigid set'' takes a slab without a side component that has
an odd number of vertices of $U$; the text takes a slab without any vertex of $U$ in a side
component, which the same count gives. The form \emph{slabs} uses \inP{Lemma~7.4} with its proof
by ladders, and this proof has no formal counterpart; the formal proof has
Lemma~\ref{lem:rails} in its place.

\section{What I examined and did not use}

\begin{itemize}[leftmargin=1.2em]
\item A slab that consists of all vertices of $M$ in its layers. Then the graph outside the slab
can have many components on each side, and a $U$-join meets the slab in a join of many
vertices. The component of the path $\pi$ and the count of side components give two vertices.
\item A proof without a lemma on two vertices. A ladder has girth four, and the two ends of one
rail form a rigid set. So some statement on rigid pairs at girth five is necessary, and a
statement of this kind needs a number such as seven: six positions exist.
\item A shorter arithmetic argument for the seven integers, by residues or by a count of the
positions in a window. I found none. The case distinction has three cases.
\item The lemma on rigid pairs inside the slab, where the path $\pi$ is given and the layers
bound the positions. The proof does not become shorter, because the condition on a shortest
path is the one that the argument uses.
\end{itemize}

\end{document}
